\section{Reading Comprehension Questions}

\noindent
{\em From Section~\ref{sec:asymptotic}}

\begin{requ}
In words, what does $f(n)=O(g(n))$ mean?
What does $f(n)=\Theta(g(n))$ mean?
\begin{requanswer}
$f(n)=O(g(n))$ means that $f(n)$ grows no faster than $g(n)$.
Another way of phrasing it is saying that $g(n)$ is an upper bound on $f(n)$.
{\em If you said that it means $g(n)$ grows faster than $f(n)$, you are not quite correct.
Go reread the definition and make sure you are clear on this point!}
$f(n)=\Theta(g(n))$ means that $f(n)$ and $g(n)$ grow at the same rate.
Another way of phrasing it is saying that $g(n)$ is a tight bound on $f(n)$.
\end{requanswer}
\end{requ}


\begin{requ}
Assume $f(n)=O(g(n))$.
\begin{lenum}
\item Does that mean $f(n)\leq g(n)$ for all values of $n$? Explain. (Hint: A complete answer should bring up two things.)
\item Does that mean that $f(n)\leq c\, g(n)$ for some constant $c$ and for all values of $n$? Explain.
\item Is it possible that $g(10)$ (for instance) is a lot smaller than $f(10)$? Explain.
\end{lenum}
\begin{requanswer}
(a) No. It means there is some $n_0$ such that whenever $n\geq n_0$, $f(n)\leq c g(n)$ (where $c>0$ is some constant).
So there are two problems with saying it means $f(n)\leq g(n)$. First, there is a constant involved.
So $f(n)\leq c\, g(n)$, not just $f(n)\leq g(n)$. Second, it does not have to hold for all values of $n$, 
but only for (i.e. $n\geq n_0$).
(b) No. Repeating what was said in the previous part, it has to hold for all values of $n$ at least as large as some
given value $n_0$, but it does not have to hold for smaller values of $n$.
(c) That is certainly possible. For instance, if $f(n)=100n$ and $g(n)=n^2$, $g(10)=100<1000=f(10)$,
but hopefully it is clear that $f(n)=O(g(n))$.
\end{requanswer}
\end{requ}

\begin{requ}
If $f(n)=O(g(n))$, does that imply that $f(n)=\Theta(g(n))$? 
If so, explain why. 
If not, give an example of functions $f$ and $g$ such that $f(n)=O(g(n))$ but $f(n)\not=\Theta(g(n))$.
\begin{requanswer}
No. For instance, if $f(n)=23n$ and $g(n)=n^2$,  then $f(n)=O(g(n))$, but  $f(n)\not=\Theta(g(n))$.
\end{requanswer}
\end{requ}

\begin{requ}
If $f(n)=\Theta(g(n))$, does that imply that $f(n)=O(g(n))$? 
If so, explain why. 
If not, give an example of functions $f$ and $g$ such that $f(n)=\Theta(g(n))$ but $f(n)\not=O(g(n))$.
\begin{requanswer}
Yes. $f(n)=\Theta(g(n))$ if and only if $f(n)=O(g(n))$ and $f(n)=\Omega(g(n))$. 
So if $f(n)=\Theta(g(n))$, then $f(n)=O(g(n))$.
\end{requanswer}
\end{requ}

\begin{requ}
Explain the difference between $f(n)=o(g(n))$ and $f(n)=O(g(n))$.
\begin{requanswer}
If $f(n)=o(g(n))$, then we know that $f(n)$ grows {\em slower than} $g(n)$. 
If $f(n)=O(g(n))$, then $f(n)$ grows {\em no faster than} $g(n)$. 
That is, $f(n)$ either grows slower than $g(n)$ or they grow at the same rate.
\end{requanswer}
\end{requ}


\begin{requ}
If you know that $f(n)=\Theta(g(n))$, does that give you more, less, or the same amount of information about 
the relationship between $f$ and $g$ than if you knew that $f(n)=O(g(n))$? Explain.
\begin{requanswer}
More. If you know that $f(n)=\Theta(g(n))$, then you know that they grow at the same rate.
But if you only know that $f(n)=O(g(n))$, it is possible that the grow at the same rate, but it is also possible
that $f(n)$ grows slower than $g(n)$.
\end{requanswer}
\end{requ}

\begin{requ}
Give two different proofs that $7n^3+4n^2-8n+27=O(n^3)$.
(Do not forget to use Theorem~\ref{thm:theta} when necessary.)
\begin{requanswer}
Proof 1: 
Notice that if $n\geq 1$, $7n^3+4n^2-8n+27\leq 7n^3+4n^2+27 \leq 7n^3+4n^3+27n^3=38n^3$.
By definition of Big-O, $7n^3+4n^2-8n+27=O(n^3)$.

\noindent
Proof 2: Notice that
$\displaystyle \lim_{n\rightarrow\infty} \frac{7n^3+4n^2-8n+27}{n^3} = 
\lim_{n\rightarrow\infty} \frac{7n^3}{n^3}+\frac{4n^2}{n^3}-\frac{8n}{n^3}+\frac{27}{n^3}
=\lim_{n\rightarrow\infty} 7+\frac{4}{n}-\frac{8}{n^2}+\frac{27}{n^3}
=7+0+0+0=7.$
Thus, $7n^3+4n^2-8n+27=\Theta(n^3)$ by Theorem~\ref{thm:limits} (part 3).
By Theorem~\ref{thm:theta}, $7n^3+4n^2-8n+27=O(n^3)$.
\end{requanswer}
\end{requ}

\begin{requ}
Prove that $3^n=o(3.1^n)$.
\begin{requanswer}
Notice that 
$\displaystyle \lim_{n\rightarrow\infty} \frac{3^n}{3.1^n} = \lim_{n\rightarrow\infty} \left(\frac{3}{3.1}\right)^n=0$
since $\frac{3}{3.1}<1$. By Theorem~\ref{thm:limits} (part 1), $3^n=o(3.1^n)$.
\end{requanswer}
\end{requ}


\noindent
{\em From Section~\ref{sec:growrates}}

\begin{requ}
Explain why $n\log n$ grows faster than $c n$ for any constant $c>0$.
That is, explain why $c n=o(n\log n)$.
Note that I am not asking for a proof of this, but
an explanation of why it makes sense.
\begin{requanswer}
Notice that both functions have a factor of $n$. Since $\log n$ grows faster than $c$ (which doesn't grow at all since
it is a constant), $n\log n$ grows faster than $c n$.

It is actually easy to prove this formally: 
$\displaystyle\lim_{n\rightarrow\infty} \frac{n\log n}{c\, n} = \lim_{n\rightarrow\infty} \frac{\log n}{c} =\infty$,
so by Theorem~\ref{thm:limits}, $n\log n=\omega(c n)$. In other words, $n\log n$ grows faster than $c n$.

\end{requanswer}
\end{requ}

\begin{requ}
We think of $\log n$ as a slow growing function.
Does that mean that given another function $f(n)$, $f(n) \log n$ grows slower than $f(n)$?
(In general, if we multiply a function by a slow growing function, does it make the function grow slower?)
Explain.
\begin{requanswer}
No! The function may grow slowly, but {\em it is still growing}. 
So you are multiplying a function by another function that is growing (even if slowly), which grows faster than 1.
Thus, $f(n)\log n$ definitely grows faster than $f(n)$.

As with the previous question, a simple limit computation gives a clear proof of this fact.
$\displaystyle\lim_{n\rightarrow\infty} \frac{f(n)\log n}{f(n)} = \lim_{n\rightarrow\infty} \frac{\log n}{1} =\infty$,
so by Theorem~\ref{thm:limits}, $f(n)\log n=\omega(f(n))$. In other words, $f(n)\log n$ grows faster than $f(n)$.

\end{requanswer}
\end{requ}

\begin{requ}
Rank the following functions in {\em increasing} order of rate of growth. Clearly indicate if two
of the functions have the same growth rate:

\medskip
$n^n$,
$7 \log_{10} n$,
$n^2+n+1$,
$7^n$,
$3n^2$,
$2^n$,
$n!$,
$\log_3 n$,
$7n\log_2 n$,
$n^3$,
$27 n$,
$8675309$, 
$n^3+n^2\log_{e} n$
\begin{requanswer}
$8675309$, 
$7 \log_{10} n$ $\sim$ $\log_3 n$,
$27 n$,
$7n\log_2 n$,
$n^2+n+1$ $\sim$ $3n^2$,
$n^3$ $\sim$ $n^3+n^2\log_{e} n$,
$2^n$,
$7^n$,
$n!$,
$n^n$
\end{requanswer}
\end{requ}

\noindent
{\em From Section~\ref{sec:algAnalysis}}

\begin{requ}
Explain why ``My algorithm only took 5 minutes to run and yours took 15 minutes, so mine is better'' is not
a complete and valid argument. (What other information is needed to make the conclusion?)
\begin{requanswer}
It is unclear what types of machines the algorithms were run on. If one was run on a cellphone and another
on a supercomputer, it is definitely not a fair comparison. 
If they were run on the same machine and on the same data, then we can compare them.
\end{requanswer}
\end{requ}

\begin{requ}
If one algorithm always takes 3 times as long to run as another algorithm, 
regardless of the size of the input,
do the two algorithms have different computational complexities? Explain.
\begin{requanswer}
No. If one always takes 3 times as long, it's just a constant multiple difference, so they have
the same growth rate. E.g. if one takes $f(n)$ time, the other takes $3f(n)$ time, and these
have the same growth rate.
Remember, when we are talking about growth rates, the constant multiples do not matter.
\end{requanswer}
\end{requ}

\begin{requ}
If algorithm $A$ has computational complexity $f(n)$, 
algorithm $B$ has computational complexity $g(n)$, and $f(n)=o(g(n))$ (that is, $g(n)$ grows faster than $f(n)$),
which algorithm is faster, $A$ or $B$? Explain.
\begin{requanswer}
Since $g(n)$ grows faster than $f(n)$, as $n$ increases, algorithm $B$ will take (relatively) longer than $A$.
In other words, algorithm $A$ is faster.
{\em Remember: Faster growth rate of computational complexity means slower algorithm!}
\end{requanswer}
\end{requ}

\begin{requ}
Someone claims to have an algorithm that can sort an array of $n$ elements in $\Theta(\log n)$ time.
Why can you be certain that they are incorrect about their algorithm?
\begin{requanswer}
You cannot even read all of the inputs in $\Theta(\log n)$ time, so this is clearly ludicrous.
It would have to take {\em at least} $\Omega(n)$ time
(and even that is not attainable, but we aren't ready for that proof yet).
\end{requanswer}
\end{requ}

\begin{requ}
Someone has an algorithm that can search an array in time $\Theta(n\log n)$. 
Is that a good or bad algorithm to solve the search problem (or is it impossible to tell)? Explain.
\begin{requanswer}
This would be an awful algorithm since searching should not take more than $\Theta(n)$ time.
\end{requanswer}
\end{requ}

\begin{requ}
Is an algorithm that can sort an array in time $\Theta(n^{1.5})$ better or worse than \verb+insertionSort+? Explain.
\begin{requanswer}
Better because \verb+insertionSort+ takes $\Theta(n^2)$ time and $n^{1.5}=o(n^2)$.
\end{requanswer}
\end{requ}

\begin{requ}
Give at least two reasons why a double-nested for loop does not always have complexity $\Theta(n^2)$.
\begin{requanswer}
One (or both) of the loops could only execute a constant number of times. 
Also, maybe the index variable(s) aren't just incremented/decremented by 1 each time through the loop.
If instead the variables are doubled or halved, for instance, the complexity might involve a logarithm or something
(e.g. maybe $\Theta(n\log n)$).
Here's a third (for free!): Maybe the outer loop executes $n$ times and the inner loop executes $m$ times.
Then the complexity is $\Theta(n\, m)$.
\end{requanswer}
\end{requ}

\begin{requ}
Algorithm $A$ has complexity $\Theta(n^2)$ and algorithm $B$ has complexity $\Theta(n\log n)$.
\begin{lenum}
\item Generally speaking, which algorithm is faster? When is it faster? Explain.
\item Are there potentially cases in which the ``worse'' algorithm is actually faster? 
If so, when might it be faster and why? If not, how do you know that it is never faster? 
\end{lenum}
\begin{requanswer}
(a) Generally speaking, $B$ is faster. Since we are given $\Theta$ bounds on the complexities, we know
the exact growth rates and $n\log n = o(n^2)$, so $B$ is definitely a faster algorithm. 
We know that $B$ is faster for large enough input. We do not know if is is faster for small inputs.
(b) As mentioned in the previous part, for small values of $n$, $A$ could possibly be faster because the
constants involved in $B$ might be much larger.
\end{requanswer}
\end{requ}

\begin{requ}
Algorithm $A$ has complexity $\Theta(n\log n)$ and algorithm $B$ has complexity $O(n^2)$.
Which algorithm is faster? Explain.
\begin{requanswer}
This is a trick question! The answer is unclear. We know that $A$ has complexity $\Theta(n\log n)$, but
we are only given a Big-O bound on the complexity of $B$. It is possible that $B$ has complexity
$\Theta(n)$ and it is faster, or that it has complexity $\Theta(n\log n)$ and it is the same as $A$,
or that it has complexity $\Theta(n^2)$ and it is slower than $A$.
\end{requanswer}
\end{requ}

\begin{requ}
If two algorithms have the same complexity, what other factors should be taken into account when choosing which
one to use? List as many as you can think of.
\begin{requanswer}
Memory usage; 
how complicated the algorithms are to implement, maintain, and debug; 
the constants; how convinced you are of the correctness of each algorithm.
\end{requanswer}
\end{requ}

